% =========================================================
% Weak (Standard) Induction
% =========================================================
\begin{tcolorbox}[colback=propbox, colframe=propborder, arc=2pt,
  left=6pt, right=6pt, top=4pt, bottom=4pt,
  title={\small\textbf{Theorem: Weak (Standard) Induction}},
  fonttitle=\small\bfseries]
Let $P(n)$ be a statement depending on $n \in \mathbb{N}$. Suppose:
\begin{enumerate}[itemsep=2pt]
  \item \textbf{Base case.} $P(n_0)$ is true for some starting value $n_0$.
  \item \textbf{Inductive step.} For all $n \geq n_0$,
    $P(n) \Rightarrow P(n+1)$.
\end{enumerate}
Then $P(n)$ is true for all $n \geq n_0$.
\end{tcolorbox}

\begin{tcolorbox}[colback=gray!4, colframe=gray!35, arc=2pt,
  left=6pt, right=6pt, top=4pt, bottom=4pt,
  title={\small\textbf{Template: Weak Induction Proof}},
  fonttitle=\small\bfseries]
Let $P(n)$ denote the statement: [\emph{write it out completely as a sentence}].

\medskip
\noindent\textbf{Base case} ($n = [n_0]$).
[\emph{Verify $P(n_0)$ directly.}]
Hence $P(n_0)$ holds.

\medskip
\noindent\textbf{Inductive step.}
Let $n \geq n_0$ and assume $P(n)$ holds. [\emph{State explicitly
what $P(n)$ says.}]
We show $P(n+1)$ holds.

[\emph{Rewrite $P(n+1)$ to expose $P(n)$ inside it.
Substitute the inductive hypothesis. Simplify.}]

Hence $P(n+1)$ holds.

\medskip
\noindent By induction, $P(n)$ holds for all $n \geq n_0$. \qed
\end{tcolorbox}

\begin{remark}[The most important discipline: write $P(n)$ as a sentence]
The single most impactful habit in induction proofs is writing the
inductive hypothesis as a \emph{complete mathematical sentence}
before the inductive step begins.

``Assume $P(n)$ holds'' is not sufficient if you have not first written
out what $P(n)$ says. Write: ``Assume
$\sum_{k=1}^n k = \frac{n(n+1)}{2}$.'' That sentence is your
inductive hypothesis. It is the precise tool you have available in
the inductive step.

Without this sentence, students typically make one of two errors:
they try to use $P(n)$ without knowing what it says, guessing at it
from context; or they prove the inductive step using information
beyond what $P(n)$ supplies, making the proof circular or invalid.
Write the sentence. It takes five seconds and eliminates the
most common induction error.
\end{remark}

\begin{remark}[The core skill of the inductive step]
Every inductive step involves the same core action: \emph{rewrite
$P(n+1)$ to expose $P(n)$ inside it, then substitute the inductive
hypothesis.}

For sum formulas, this means: split off the last term.
\[
\sum_{k=1}^{n+1} k = \left(\sum_{k=1}^n k\right) + (n+1).
\]
Now substitute $P(n)$: replace $\sum_{k=1}^n k$ with $\frac{n(n+1)}{2}$.
Simplify to get the $P(n+1)$ form.

For divisibility, this means: write $4^{n+1} - 1 = 4(4^n - 1) + 3$
to expose $4^n - 1$ inside $4^{n+1} - 1$.

In both cases, the step is: factor out the $P(n)$ term, substitute,
and verify the arithmetic produces $P(n+1)$.

When the rewriting step is not obvious, it usually means $P(n)$ needs
to be strengthened (see Section~\ref{sec:choosing-pn}).
\end{remark}

\begin{example}[Sum formula]
\textbf{Proposition.}
For all $n \geq 1$:
$\displaystyle\sum_{k=1}^n k = \frac{n(n+1)}{2}$.

\medskip
\noindent\textit{Proof.}
Let $P(n)$ denote: $\displaystyle\sum_{k=1}^n k = \frac{n(n+1)}{2}$.

\medskip
\noindent\textbf{Base case} ($n=1$).
$\sum_{k=1}^1 k = 1 = \frac{1 \cdot 2}{2}$.
Hence $P(1)$ holds.

\medskip
\noindent\textbf{Inductive step.}
Let $n \geq 1$ and assume
$\displaystyle\sum_{k=1}^n k = \frac{n(n+1)}{2}$. \quad
\textit{(Inductive hypothesis.)}

We show $\displaystyle\sum_{k=1}^{n+1} k = \frac{(n+1)(n+2)}{2}$.

Split off the last term:
\begin{align*}
\sum_{k=1}^{n+1} k
&= \left(\sum_{k=1}^n k\right) + (n+1) \\
&= \frac{n(n+1)}{2} + (n+1)
  \tag{inductive hypothesis} \\
&= (n+1)\left(\frac{n}{2} + 1\right)
  = \frac{(n+1)(n+2)}{2}.
\end{align*}

Hence $P(n+1)$ holds. By induction, $P(n)$ holds for all $n \geq 1$. \qed
\end{example}

\begin{example}[Divisibility]
\textbf{Proposition.}
For all $n \geq 0$, $3 \mid (4^n - 1)$.

\medskip
\noindent\textit{Proof.}
Let $P(n)$ denote: $3 \mid (4^n - 1)$.

\medskip
\noindent\textbf{Base case} ($n=0$).
$4^0 - 1 = 0 = 3 \cdot 0$. Hence $3 \mid 0$ and $P(0)$ holds.

\medskip
\noindent\textbf{Inductive step.}
Assume $3 \mid (4^n - 1)$, i.e., $4^n - 1 = 3m$ for some $m \in \mathbb{Z}$.

Expose $4^n - 1$ inside $4^{n+1} - 1$:
\[
4^{n+1} - 1 = 4 \cdot 4^n - 1
= 4(4^n - 1) + 4 - 1
= 4(4^n - 1) + 3
= 4 \cdot 3m + 3
= 3(4m + 1).
\]
Hence $3 \mid (4^{n+1}-1)$ and $P(n+1)$ holds.

By induction, $P(n)$ holds for all $n \geq 0$. \qed

\medskip
\noindent
The rewriting step $4^{n+1} - 1 = 4(4^n - 1) + 3$ is the key move:
it exposes the inductive hypothesis $3 \mid (4^n - 1)$ inside
the goal $3 \mid (4^{n+1} - 1)$.
\end{example}
